\begin{center}
    {\LARGE \textbf{2004无机化学}} \\[2ex]
    {\large 时量180分钟 \quad 满分150分}
\end{center}

\vspace{0.5cm}
\noindent\textbf{注意事项:}
\begin{enumerate}[label=\arabic*., parsep=0pt, itemsep=0pt]
    \item 答题前,考生先将自己的姓名、学号、考场号、座位号填写在答题卡上,将条形码准确粘贴在考生信息条形码粘贴区。
    \item 考试结束后,将本试卷与答题纸一并收回。
    \item 回答各个问题时,将答案写在答题纸上,写在本试卷上无效。
\end{enumerate}

\vspace{1cm}
\begin{center}
    {\Large \textbf{理论部分(共70分)}}
\end{center}

\begin{enumerate}[label=\arabic*., start=1, itemsep=1.5ex]
    \item 依次写出第五周期元素的名称和元素符号。在元素符号下划横线以表示出核外有10个d电子的元素；在元素名称下划横线以表示出核外5s轨道中只有1个电子的元素。(18 pts.)

    \item 用价层电子对互斥理论判断\ce{NO2}分子的构型。再用杂化轨道理论分析\ce{NO2}分子的成键情况，并对存在的离域π键给出理论上的支持或实验事实上的支持。(14 pts.)

    \item \ce{[Co(NH3)6]Cl3}是一种逆磁性的橙黄色晶体，试写出该化合物的名称。用价键理论分析其内界的成键情况。画简图表示中心d轨道在分裂后的d轨道中的排布方式，并计算其晶体场稳定化能。(14 pts.)

    \item 已知：

        $E^\ominus_A(\ce{H3AsO4/H3AsO3}) = 0.56\ \text{V}$

        $E^\ominus_A(\ce{I2/I^-}) = 0.54\ \text{V}$

        试写出这两个电对所对应的电极反应式；分别计算出这两个电对在$\text{pH}=7.0$时的$E$值；根据计算结果写出这两个电对中的四种物质均以标准态浓度共存于$\text{pH}=7.0$的水溶液中时的反应方程式。(12 pts.)

    \item 利用元素电势图，举例解答下面问题。(12 pts.)
    
    $E^\ominus/\text{V} \quad 
    \ce{Cr^{3+}} 
    \xline{\raisebox{0.5ex}{$-0.41$}} 
    \ce{Cr^{2+}} 
    \xline{\raisebox{0.5ex}{$-0.86$}} 
    \ce{Cr}$

    \begin{enumerate}[label=(\arabic*), itemsep=1ex]
        \item 同一个化学反应是否可以设计成几个不同的原电池？用电池符号表示出来。
        \item 这几个原电池的电动势$E^\ominus$可否不相等？用计算结果表示出来。
        \item 根据$E^\ominus$计算出的电池反应的$\Delta_r G_m^\ominus$是否相同？用计算结果说明问题。
    \end{enumerate}

\end{enumerate}

\vspace{1cm}
\begin{center}
    {\Large \textbf{元素部分(共80分)}}
\end{center}

\begin{enumerate}[label=\arabic*., start=1, itemsep=1.5ex]
    \item 设计不使用\ce{H2S}的方案，分离下列离子：\ce{Ag^{+}},\ce{Hg^{2+}},\ce{Cd^{2+}},\ce{Zn^{2+}},\ce{Al^{3+}},\ce{Fe^{3+}} (10 pts.)

    \item 用化学方程式完成由\ce{CoSO4\cdot 7H2O}制备无水\ce{CoCl2}，写出应有的实验现象。(10 pts.)

    \item 用化学方程式表示下列转化过程：\ce{S -> SO2 -> Na2SO3 -> Na2S2O3 -> Na2S4O6} (10 pts.)
    
    \item 三个试管中分别盛有\ce{MnO2}、\ce{PbO2}和\ce{Fe3O4}，试设计出鉴别方案。 (10 pts.)
    
    \item 金属\ce{M}的硝酸盐\ce{A}为无色晶体，\ce{A}与水反应得白色沉淀物\ce{B}和溶液\ce{C}。向\ce{C}中通入\ce{H2S}，生成黑色沉淀物\ce{D}。\ce{D}不溶于\ce{NaOH}溶液，但可溶于盐酸中。向\ce{C}中滴加\ce{NaOH}溶液，生成白色沉淀物\ce{E}，\ce{E}不溶于过量\ce{NaOH}溶液。向氯化亚锡的强碱性溶液中加入\ce{C}，将有金属\ce{M}生成。试给出\ce{M}、\ce{A}、\ce{B}、\ce{D}和\ce{E}的化学式。(10 pts.)
    
    \item 将绿色固体\ce{A}溶于水后，通入足量的\ce{SO2}得溶液\ce{B}。向\ce{B}中加入\ce{NaOH}溶液生成白色沉淀物\ce{C}，\ce{C}不溶于过量的\ce{NaOH}溶液，再加入\ce{H2O2}，则沉淀物变为暗棕色的\ce{D}。将\ce{D}滤出与\ce{KClO3}、\ce{KOH}共熔可制得\ce{A}。试给出\ce{A}、\ce{B}、\ce{C}、\ce{D}的化学式及由\ce{D}和\ce{KClO3}、\ce{KOH}共熔制\ce{A}的反应方程式。(12 pts.)
    
    \item 写出八种金属卤化物难溶盐的化学式，并指明各自的颜色。(8 pts.)
    
    \item 试通过Li和Mg、Be和Al、B和Si的相似性比较，从事实上说明"对角线规则"的存在，并以理论说明"对角线规则"成立的根据。(10 pts.)
\end{enumerate}


